\chapter{Quantized cotangent charts and their central quotient}
\label{chap:qc-quantum-collision}

A constraint on coordinates can be imposed before or after passing to local charts.
For differential operators, the two procedures must retain the order of multiplication.
That order produces a scalar correction which is absent from the equation of the classical constraint.
The simplest projective example makes all three features visible: the constraint, the coordinate changes, and the correction.

Let $k$ be a field of characteristic zero.
Describe the projective line $\mathbb P^1$ by two affine charts, with polynomial coordinate algebras $k[x]$ and $k[x']$.
Their intersection has Laurent coordinate algebra $k[x,x^{-1}]$, and $x'=x^{-1}$.
Its $k$-valued points are nonzero pairs $(a_1,a_2)$ modulo simultaneous nonzero scaling.
The chart coordinates are $x=a_2/a_1$ and $x'=a_1/a_2$, respectively.
A cotangent vector in the first chart is written $v\,dx$.
The equality $v\,dx=v'\,dx'$ gives
\[
 x'=x^{-1},\qquad v'=-x^2v.
\]
Thus the cotangent bundle is obtained by gluing two polynomial planes along these formulas.
The zero cotangent vectors form a projective line inside this surface.

There is also a description by a constraint.
Take a column $a=(a_1,a_2)$ and a row $p=(p_1,p_2)$, and put $\mu=p_1a_1+p_2a_2$.
Impose $a\ne0$ and $\mu=0$, then identify $(a,p)$ with $(c a,c^{-1}p)$.
For the line $L=ka$, the rule $\bar u\mapsto p(u)a$ defines a map $k^2/L\to L$.
It is well-defined because $p(a)=0$, and scaling does not change it.
Such maps are precisely the cotangent vectors: the tangent space is $\operatorname{Hom}(L,k^2/L)$, and composition followed by trace pairs the two spaces.
Choosing $a$ reconstructs $p$ from the map, uniquely up to the displayed scaling.
This proves the cotangent description.
On $a_1\ne0$ its coordinates are $x=a_2/a_1$ and $v=a_1p_2$.
Indeed, the one-form $p_1\,da_1+p_2\,da_2$ restricts at $\mu=0$ to $v\,dx$.

Quantization replaces commuting momenta and coordinates by prescribed commutators.
The constraint becomes an operator, and its chosen value enters the coordinate transition.

\section{The moment operator and an ordered basis}

Let $R=k[\hbar,\nu]$.
Both parameters are central and algebraically independent.
The algebra $\mathsf W$ is generated over $R$ by $a_1,a_2,p_1,p_2$, with
\[
 [p_i,a_j]=\hbar\delta_{ij},\qquad [a_i,a_j]=[p_i,p_j]=0,
 \qquad [u,v]=uv-vu.
\]
All elements are finite linear combinations of words.
Every tensor product and polynomial algebra in this chapter is algebraic.
No inverse of $\hbar$ or completion in $\hbar$ is taken.

\begin{lemma}\label{lem:qc-weyl-basis}
The monomials $a_1^{i_1}a_2^{i_2}p_1^{j_1}p_2^{j_2}$ form an $R$-basis of $\mathsf W$.
\end{lemma}
\begin{proof}
Moving each $p_i$ to the right uses the defining relation and produces shorter correction terms.
Ordering first by length and then by the number of reversed pairs proves that this procedure terminates.
Thus the displayed monomials span.
For independence, act on $R[a_1,a_2,t_1,t_2]$ by multiplication with $a_i$ and by
\[
 p_i=\hbar\frac{\partial}{\partial a_i}+t_i.
\]
These operators satisfy every relation.
Applying a normally ordered linear combination to $1$ replaces each $p_i$ by $t_i$.
Independence of polynomial monomials proves independence of the proposed basis.
This argument remains valid at $\hbar=0$.
\end{proof}

Give $a_i$ weight $1$ and $p_i$ weight $-1$.
The relations preserve weight, so $\mathsf W$ is a direct sum of its weight spaces.
Its weight-zero subalgebra is denoted $\mathsf W^0$.
Retaining weight zero is the algebraic scaling-invariant operation: scaling by $c$ multiplies a word of weight $m$ by $c^m$.
For maps preserving weight, kernels and images decompose by weight.
Consequently retaining weight zero preserves exact sequences.

Define the moment operator
\[
 M=a_1p_1+a_2p_2.
\]
The commutator product identity gives $[M,a_i]=\hbar a_i$ and $[M,p_i]=-\hbar p_i$.
Hence $[M,w]=m\hbar w$ for a word of weight $m$.
In particular, $M$ is central in $\mathsf W^0$.
The reduced algebra at period $\nu$ is
\begin{equation}\label{eq:qc-reduction}
 \mathsf A_{\hbar,\nu}=\mathsf W^0/(M-\hbar\nu).
\end{equation}
The quotient is two-sided because its defining element is central.
It is also the weight-zero part of the quotient by the left ideal $\mathsf W(M-\hbar\nu)$.
Indeed, the weight-zero part of that ideal is $\mathsf W^0(M-\hbar\nu)$.

The parameter $\nu$ specifies the value of a quantum moment operator.
It is independent of any extra normal coordinate introduced to express a potential.

\section{The relative bar retains the moment relation}

The quotient in \eqref{eq:qc-reduction} forgets how the relation is imposed.
A two-term complex retains that information and admits a direct comparison with an ordered bar.
The comparison works for any central element, before its injectivity is considered.

Let $A$ be a unital associative algebra over a commutative ring $B$, and let $w$ be central in $A$.
Let $B[t]$ act on $A$ by $t\mapsto w$, and on $B$ by $t\mapsto0$.
Put $I=tB[t]$.
The normalized two-sided bar has degree $-n$ term
\[
 \mathcal B^{-n}=A\otimes_B I^{\otimes_B n},\qquad n\ge0.
\]
Write $a[m_1|\cdots|m_n]$ for the element with letters $t^{m_i}$, where every $m_i\ge1$.
The empty word is an element of $A$.
Define a differential of degree $1$ by
\begin{align*}
 d\bigl(a[m_1|\cdots|m_n]\bigr)
 &=aw^{m_1}[m_2|\cdots|m_n]\\
 &\quad+\sum_{i=1}^{n-1}(-1)^i
 a[m_1|\cdots|m_i+m_{i+1}|\cdots|m_n].
\end{align*}
There is no final term because every letter has augmentation zero.
Each pair of successive multiplications occurs twice in $d^2$, with opposite signs.
The initial pair cancels because $w^{m_1}w^{m_2}=w^{m_1+m_2}$.
Thus $d^2=0$.

A shuffle interleaves two words while retaining the order inside each word.
Give each letter odd degree and attach the sign of the interleaving permutation.
The product on $\mathcal B$ multiplies the two coefficients in their given order and sums these signed shuffles.
The empty word with coefficient $1$ is its unit.
An interleaving of three ordered lists is independent of the order in which two lists are first interleaved.
This proves associativity with the same signs.
For the differential, faces joining letters from different lists cancel in pairs.
Faces within one list give the corresponding factor differential.
The initial face obeys the same rule because each $w^m$ is central in $A$.
The product therefore satisfies the graded Leibniz identity.

Let
\[
 C=A\otimes_B\Lambda(\epsilon),\qquad |\epsilon|=-1,
 \qquad d\epsilon=w.
\]
Here $\epsilon$ commutes with $A$ and satisfies $\epsilon^2=0$.
The differential respects this relation since $w\epsilon-\epsilon w=0$.

\begin{proposition}\label{prop:qc-bar-contraction}
There are unital differential-algebra maps
\[
 i:C\longrightarrow\mathcal B,\qquad
 \pi:\mathcal B\longrightarrow C,
\]
given by $i(a)=a$, $i(a\epsilon)=a[1]$, and
\[
 \pi(a)=a,\qquad
 \pi(a[m])=aw^{m-1}\epsilon,\qquad
 \pi(a[m_1|\cdots|m_n])=0\quad(n\ge2).
\]
They satisfy $\pi i=1$ and $dh+hd=1-i\pi$, where $h$ has degree $-1$ and
\[
 h\bigl(a[m_1|\cdots|m_n]\bigr)=
 \begin{cases}
 (-1)^n a[m_1|\cdots|m_{n-1}|m_n-1|1],&m_n\ge2,\\
 0,&m_n=1,
 \end{cases}
 \qquad h(a)=0.
\]
No regularity assumption on $w$ is needed.
\end{proposition}
\begin{proof}
The square of $[1]$ is $[1|1]-[1|1]=0$.
It commutes with every degree-zero coefficient, so $i$ preserves products and the differential.
The formula for $\pi$ commutes with the differential on one-letter words.
On a two-letter boundary the two possible images are
$aw^{m_1+m_2-1}\epsilon$ and its negative.
All longer boundaries map to zero.
Products of two positive-length words have length at least two, while the product of their images contains $\epsilon^2$ or is zero.
Multiplication with a degree-zero coefficient is preserved by centrality of $w$.
Thus $\pi$ is also a differential-algebra map.

For one letter with $m\ge2$,
\[
 dh(a[m])=a[m]-aw^{m-1}[1].
\]
For $m=1$, both sides of $dh+hd=1-i\pi$ vanish.
For a word of length at least two with last exponent at least two, every face before the last pair cancels between $dh$ and $hd$.
The face merging the final two letters of $h$ returns the original word with sign $(-1)^n(-1)^n=1$.
The remaining face cancels the term obtained by first merging the last pair of the original word.
If the last exponent is one, $h$ vanishes on the original word.
In $hd$ only that last merged face survives, and its two signs again multiply to $1$.
This proves the homotopy identity in every length.
The other assertions follow directly from the formulas.
\end{proof}

Taking $B=R$, $A=\mathsf W^0$, and $w=M-\hbar\nu$ gives the relative moment bar.
It computes the derived tensor product over $R[z]$, where $z$ acts as $M$ on $\mathsf W^0$ and as $\hbar\nu$ on $R$.
Here the derived tensor product means tensoring with a free resolution before taking cohomology.
The two-term resolution of $R[z]/(z-\hbar\nu)$ is free and exact because multiplication by this monic polynomial is injective.
It gives exactly $C$ after tensoring with $\mathsf W^0$.

If $A=k[u]/(u^2)$ and $w=u$, the element $u\epsilon$ is a nonzero cycle in degree $-1$ of $C$.
There is no term of degree $-2$, so this cycle cannot be a boundary.
Proposition~\ref{prop:qc-bar-contraction} retains this class in the relative bar.
Replacing either complex by the ordinary quotient $A/(u)$ would lose it.

For these coefficients, $w$ is a non-zero-divisor.
Filter $\mathsf W$ by total word length.
Its associated graded is the polynomial ring over the domain $R$, by Lemma~\ref{lem:qc-weyl-basis}.
Leading terms of nonzero products cannot vanish, so $\mathsf W$ and its subalgebra $\mathsf W^0$ are domains.
The ordered basis also shows $M-\hbar\nu\ne0$.
Thus $C$ has cohomology only in degree zero, equal to \eqref{eq:qc-reduction}.

Changing the invariant operation changes this conclusion.
For the one-dimensional Lie algebra generating scaling, the cochain complex on a weight-$m$ space has two terms and differential multiplication by $m$.
Its second term has an odd degree-one generator $c$.
For $m\ne0$, division by $m$ contracts the complex.
At weight zero its differential vanishes and both terms remain.
Applied after the moment quotient, it leaves
$\mathsf A_{\hbar,\nu}\otimes\Lambda(c)$ in cohomology.
This additional class explains the difference between Lie-algebra cochains and retaining algebraic scaling invariants.

\section{Operator charts and their overlap}

On $a_1\ne0$, write
\[
 z=a_1,\qquad x=a_2/a_1,\qquad v=a_1p_2.
\]
The localized Weyl algebra consists of ordered differential expressions with coefficients in $R[a_1,a_1^{-1},a_2]$.
Differentiation of these Laurent coefficients extends its product.
Its presentation in the new coordinates has
\[
 [v,x]=\hbar,\qquad [M,z]=\hbar z,
\]
and all other commutators between $x,v,z,M$ vanish.
The inverse assignments are
\[
 a_1=z,\quad a_2=zx,\quad
 p_2=z^{-1}v,\quad p_1=z^{-1}(M-xv).
\]
They satisfy the original relations.
For example, both products $p_1p_2$ and $p_2p_1$ become
$z^{-2}(M-\hbar-xv)v$.
The remaining commutators follow from the two displayed relations.
These inverse assignments prove the presentations agree.
Only $z$ has nonzero weight.
An ordered basis in powers of $z,x,v,M$ therefore identifies the weight-zero algebra with the Weyl chart in $x,v$ and the central polynomial variable $M$.

On $a_2\ne0$ use $x'=a_1/a_2$ and $v'=a_2p_1$.
Before the moment relation, the transition is
\[
 x'=x^{-1},\qquad v'=xM-x^2v.
\]
After the relation it becomes
\begin{equation}\label{eq:qc-operator-transition}
 x'=x^{-1},\qquad v'=-x^2v+\hbar\nu x.
\end{equation}
For a Laurent polynomial $f$, the product rule is
$vf=fv+\hbar f'$.
It follows for positive powers by induction, for $x^{-1}$ by differentiating the inverse relation, and then for every Laurent polynomial by linearity.
It gives $[v',x']=\hbar$ and the inverse formula
$v=-x'^2v'+\hbar\nu x'$.
Thus \eqref{eq:qc-operator-transition} is an algebra isomorphism on the overlap.

Denote the resulting sheaf of algebras on $\mathbb P^1$ by $\mathcal D_{\hbar,\nu}$.
A sheaf here assigns compatible local expressions to open sets, with unique gluing whenever the expressions agree on intersections.
The two chart presentations and their restrictions specify these assignments.
At $\hbar=0$, the generators commute and the transition is the cotangent transition.
At $\hbar=1$, the local generator $v$ acts as $\partial_x$.
For an integer period $\nu=n$, apply $a_2\partial_{a_1}$ to a homogeneous expression $a_1^n f(x)$.
The result is $a_1^n(-x^2 f'(x)+nx f(x))$.
This proves the transition for differential operators acting on the line bundle whose local frames satisfy $a_2^n=x^n a_1^n$.
That bundle is denoted $\mathcal O(n)$.
For arbitrary $\nu$, the same algebraic transition defines twisted differential operators without requiring such a line bundle.

\section{Three global operators}

Set
\begin{equation}\label{eq:qc-three-operators}
 E=v,\qquad F=-x^2v+\hbar\nu x,\qquad H=\hbar\nu-2xv.
\end{equation}
On the second chart these become, respectively,
$-x'^2v'+\hbar\nu x'$, $v'$, and $2x'v'-\hbar\nu$.
They therefore extend to global sections.
The product rule gives
\[
 [H,E]=2\hbar E,\qquad [H,F]=-2\hbar F,
 \qquad [E,F]=\hbar H.
\]
Let $U_\hbar$ be the associative algebra over $R$ generated by $E,F,H$ with these relations.
Thus the three global sections define an algebra map from $U_\hbar$.

\begin{lemma}\label{lem:qc-enveloping-basis}
The words $F^iH^jE^\ell$ form an $R$-basis of $U_\hbar$.
The element
\[
 \Omega=H^2+2\hbar H+4FE
\]
is central.
\end{lemma}
\begin{proof}
Use the rules $EH=HE-2\hbar E$, $HF=FH-2\hbar F$, and $EF=FE+\hbar H$.
Word length and then the number of inversions decrease under each reduction.
Disjoint reductions commute.
The only overlapping pair of reductions is the word $EHF$.
Both orders reduce it to
\[
 FHE-4\hbar FE+\hbar H^2-2\hbar^2H.
\]
Induction on the reduction order therefore makes every normal form unique.
The normal-form map kills the defining relations and their multiples, so the normal words are independent in the quotient.

The commutator with $H$ vanishes because $[H,FE]=0$.
For $E$, the commutator of $\Omega$ is
\[
 2\hbar(HE+EH)+4\hbar^2E-4\hbar HE=0.
\]
For $F$, it is
$-2\hbar(HF+FH)-4\hbar^2F+4\hbar FH=0$.
These equations use exactly the defining relations.
\end{proof}

Inside $\mathsf W^0$, put $X_{ij}=a_ip_j$.
Then $E=X_{12}$, $F=X_{21}$, and $H=X_{11}-X_{22}$.
The two diagonal elements commute and
$FE=X_{11}X_{22}+\hbar X_{22}$.
Consequently
\begin{equation}\label{eq:qc-casimir}
 \Omega=(X_{11}+X_{22})^2+2\hbar(X_{11}+X_{22})
 =M(M+2\hbar).
\end{equation}
Imposing the moment relation gives $\Omega=\hbar^2\nu(\nu+2)$.
The term $2\hbar^2\nu$ comes from multiplying a momentum past a coordinate.

\section{The global algebra from its symbols}

An increasing filtration $F_nA$ is a sequence of submodules with
$F_nA\,F_mA\subset F_{n+m}A$ and $\bigcup_nF_nA=A$.
Its associated graded algebra is $\bigoplus_nF_nA/F_{n-1}A$.
The class of an element of highest order $n$ in this quotient is its leading symbol.
Give momenta differential order one and coordinates order zero.
All parameters in $R$ have order zero.

\begin{lemma}\label{lem:qc-invariant-symbols}
The associated graded algebra of $\mathsf W^0$ is
\[
 R[q_{11},q_{12},q_{21},q_{22}]/(q_{11}q_{22}-q_{12}q_{21}),
 \qquad q_{ij}=a_i p_j.
\]
The four operators $X_{ij}$ generate $\mathsf W^0$.
\end{lemma}
\begin{proof}
The ordered basis identifies the symbols with commuting monomials having equal total coordinate and momentum degrees.
Pair each coordinate factor with a momentum factor to express such a monomial in the $q_{ij}$.
The determinant relation holds for their images.
Replace every occurrence of $q_{11}q_{22}$ by $q_{12}q_{21}$.
The resulting monomials have exponents $(a,b,c,d)$ with $\min(a,d)=0$.
Their images have exponents $(a+b,c+d,a+c,b+d)$ in $(a_1,a_2,p_1,p_2)$.
Two exponent quadruples with the same image differ by $(t,-t,-t,t)$.
The condition $\min(a,d)=0$ determines $t$ uniquely.
Thus the reduced monomials have distinct images, proving the asserted presentation.

The same pairing of factors in the noncommutative algebra gives a product of the $X_{ij}$ with the required leading monomial.
Its other terms have smaller differential order.
Subtract them and use induction on order.
Every element has finite order, so this proves generation.
\end{proof}

The leading symbol of $M-\hbar\nu$ is $q_{11}+q_{22}$.
The symbol ring in Lemma~\ref{lem:qc-invariant-symbols} is a subring of a polynomial domain, so this nonzero symbol is a non-zero-divisor.
For a central element with regular leading symbol, the associated graded of its quotient is the quotient by that symbol.
Indeed, its ideal consists of its multiples, and the order of each nonzero multiple is the sum of the two orders.
Hence the leading symbols of the ideal are exactly the multiples of its leading symbol.
It follows that
\[
 \operatorname{gr}\mathsf A_{\hbar,\nu}
 =R[e,f,\eta]/(\eta^2+4ef),
 \qquad e=q_{12},\quad f=q_{21},\quad \eta=q_{11}-q_{22}.
\]
Here division by two expresses the diagonal symbols in terms of their trace and difference.

The same argument applies to
\[
 Q=U_\hbar/(\Omega-\hbar^2\nu(\nu+2)).
\]
Its leading relation is $\eta^2+4ef$, which is monic in $\eta$ and therefore a non-zero-divisor.
Thus $\operatorname{gr}Q$ is the same quadratic ring.
The map $Q\to\mathsf A_{\hbar,\nu}$ is surjective because the $X_{ij}$ generate and their trace has been fixed.
Its map on symbols is the identity.
An element of its kernel would have a nonzero leading symbol in that kernel, which is impossible.
Therefore $Q\cong\mathsf A_{\hbar,\nu}$.

To compare with all global operators, the remaining question is whether an arbitrary compatible pair of chart expressions is generated by $E,F,H$.
This requires an overlap calculation.

\begin{lemma}\label{lem:qc-laurent-splitting}
Let $\mathcal D_{\le n}$ denote operators of order at most $n$ in $\mathcal D_{\hbar,\nu}$.
Every order-at-most-$n$ expression on the overlap is a difference of expressions from the two charts.
The leading symbols of global operators of order $n$ are exactly
\[
 \sum_{j=0}^{2n}c_j x^j v^n,\qquad c_j\in R.
\]
Every such symbol lifts to a global operator.
\end{lemma}
\begin{proof}
On symbols, the transition is $(v')^n=(-1)^n x^{2n}v^n$.
For the overlap difference map, the leading coefficient therefore has the form
\[
 a(x)-(-1)^n x^{2n}b(x^{-1}),\qquad
 a\in R[x],\quad b\in R[x'].
\]
The first summand contains all powers of exponent at least zero.
The second contains all powers of exponent at most $2n$.
Since $n\ge0$, together they contain every Laurent monomial.
Given an overlap operator, split its leading coefficient in this way.
Subtract the corresponding chart expressions and repeat at lower order.
This finite induction proves the first assertion.

For a compatible leading symbol, $a(x)=(-1)^n x^{2n}b(x^{-1})$.
Both sides are regular on their charts precisely when $a$ has degree at most $2n$.
Choose chart operators with these leading coefficients.
Their discrepancy has order at most $n-1$.
The first assertion at that order expresses the discrepancy as a chart difference.
Subtract these lower-order corrections to obtain a global lift.
For $n=0$, no correction is needed.
\end{proof}

The kernel and cokernel of the overlap difference form the degree-zero and degree-one cohomology of the two-chart complex.
The lemma computes this complex directly: its degree-one cohomology vanishes at every finite order.
The leading global symbols form the ring
\[
 \bigoplus_{n\ge0}\operatorname{span}_R\{x^jv^n:0\le j\le2n\}.
\]
It is generated by $v,-x^2v,-2xv$, the symbols of $E,F,H$.
Its only relation is $\eta^2+4ef=0$.
For an explicit verification, the monomials $f^i\eta^j e^k$ with $j=0$ or $1$ map to nonzero scalar multiples of
$x^{2i+j}v^{i+j+k}$.
At a fixed total degree $n$, these give once each the exponents $0,1,\ldots,2n$.
They are therefore a basis, and no further relation can hold.

\begin{theorem}\label{thm:qc-global-algebra}
The displayed maps give isomorphisms of filtered $R$-algebras
\[
 \mathsf A_{\hbar,\nu}
 \cong U_\hbar/(\Omega-\hbar^2\nu(\nu+2))
 \cong\Gamma(\mathbb P^1,\mathcal D_{\hbar,\nu}).
\]
Each algebra has $R$-basis
\[
 \{F^iH^jE^\ell:i,\ell\ge0,\ j\in\{0,1\}\}.
\]
Its order-at-most-$n$ part is free of rank $(n+1)^2$.
The presentations and comparison with global sections commute with every homomorphism from $R$ to a commutative $k$-algebra.
\end{theorem}
\begin{proof}
The first isomorphism has been proved above.
The map from its central quotient to global sections is an isomorphism on leading symbols by Lemma~\ref{lem:qc-laurent-splitting}.
Subtracting a lift of the leading symbol proves surjectivity by decreasing order.
The same leading-symbol argument proves injectivity.

The specified monomials lift the basis of each graded piece.
They span by decreasing order and are independent by taking the largest order in any finite relation.
There are $2m+1$ basis elements of exact order $m$.
Summing for $0\le m\le n$ gives $(n+1)^2$.

All reductions have coefficients in $R$, with no division by either parameter.
The basis remains a basis after any coefficient homomorphism.
The weight-zero ordered basis commutes with coefficient change, and the quotient presentation commutes with tensor product.
The Laurent splitting in Lemma~\ref{lem:qc-laurent-splitting} also uses no such division and works over the new coefficient ring, including rings with zero divisors.
Its symbol basis and the same decreasing-order arguments prove the global comparison after coefficient change.
\end{proof}

The freeness proves flatness: tensoring a free module preserves exact sequences.
The moment sequence $0\to\mathsf W^0\xrightarrow{M-\hbar\nu}\mathsf W^0\to\mathsf A_{\hbar,\nu}\to0$ splits as a sequence of $R$-modules.
Choose a preimage of each basis element of its free quotient to obtain a splitting.
It therefore remains exact after every coefficient change, including a nonflat one.
The moment Koszul complex continues to resolve the ordinary quotient on those fibres.
At $\hbar=0$ the algebra becomes $k[\nu,e,f,\eta]/(\eta^2+4ef)$.
The surface defined by this equation is a quadratic cone, independent of $\nu$.
The stable cotangent surface retains additional geometry.
Its zero section is an entire projective line, whereas all three global fibre-linear functions vanish there.
Thus passing to the global algebra contracts this zero section to the vertex of the cone.

\section{Normal elimination as a multiplicative global comparison}

To the Weyl algebra in $x,v$, adjoin central polynomial variables $\mu,\lambda$ and use
\[
 x'=x^{-1},\qquad v'=x\mu-x^2v.
\]
The variables $\mu,\lambda$ agree on the two charts.
Put $s=\mu-\hbar\nu$.
The expression $\lambda s$ is a central potential.
On a free right module with one even and one odd basis vector, define
\[
 D=\begin{pmatrix}0&s\\\lambda&0\end{pmatrix},\qquad
 L=\begin{pmatrix}0&0\\1&0\end{pmatrix}.
\]
The identity $D^2=\lambda s I$ makes this a matrix factorization of that potential.
Its endomorphisms form a parity differential algebra $\mathcal E$, with
$\delta T=DT-(-1)^{|T|}TD$.
Here parity is degree modulo two, and even and odd matrices are diagonal and off-diagonal, respectively.
Centrality of $D^2$ gives $\delta^2=0$.

Let $\widetilde{\mathcal D}$ be the unreduced chart sheaf with variable $\mu$ and without $\lambda$.
Set
\[
 \mathcal K=\widetilde{\mathcal D}\otimes\Lambda(\epsilon),
 \qquad d\epsilon=s.
\]
For this comparison, reduce the cohomological grading of the two-term complex to parity.
The extra generator $\epsilon$ is fixed on overlaps.

\begin{proposition}\label{prop:qc-normal-sheaf}
There is a strict diagram of unital parity differential-algebra sheaves
\[
 \mathcal D_{\hbar,\nu}\ \xleftarrow{\ p\ }\
 \mathcal K\ \xrightarrow{\ i\ }\ \mathcal E.
\]
Both arrows are quasi-isomorphisms on each chart and its overlap.
The first sets $s=\epsilon=0$.
The second sends coefficients to scalar matrices and $\epsilon$ to $L$.
\end{proposition}
\begin{proof}
For an arbitrary unital coefficient algebra $A$, central polynomial variables $s,\lambda$ act injectively on $A[s,\lambda]$.
Direct multiplication gives
\[
 \delta\operatorname{diag}(a,b)=
 \begin{pmatrix}0&s(b-a)\\\lambda(a-b)&0\end{pmatrix},
 \qquad
 \delta\begin{pmatrix}0&c\\d&0\end{pmatrix}=(sd+\lambda c)I.
\]
An even cycle is scalar, and its boundaries are precisely scalar matrices with coefficients in $(s,\lambda)$.
For an odd cycle, reduce $sd+\lambda c=0$ modulo $s$.
Injectivity of multiplication by $\lambda$ implies $c=sr$.
Injectivity of multiplication by $s$ then gives $d=-\lambda r$.
This matrix is the boundary of $\operatorname{diag}(0,r)$.
The endomorphism cohomology is therefore $A$ in even parity and zero in odd parity.
The two-term complex $A[s]\otimes\Lambda(\epsilon)$ has the same cohomology.

The identities $L^2=0$ and $\delta L=sI$ prove that $i$ preserves multiplication and differentials.
Scalar matrices commute with $L$, although their coefficients need not commute with each other.
Both arrows induce the identity on the computed even cohomology.
Their formulas commute strictly with the unreduced coordinate transition, which fixes $s,\lambda,L$.
This proves the sheaf assertion and every local quasi-isomorphism.
\end{proof}

Local inverse inclusions have an explicit multiplicative homotopy.
Let $t$ have degree zero and let $dt$ have degree one, with $(dt)^2=0$.
Write differential forms to the left of matrix coefficients.
The total differential satisfies
$d(\alpha T)=d\alpha\,T+(-1)^{|\alpha|}\alpha\delta T$.
On the overlap, send
\[
 x'\longmapsto x^{-1}I,\qquad
 v'\longmapsto(-x^2v+\hbar\nu x+t xs)I+dt\,xL.
\]
The differential of $t xsI$ is $dt\,xsI$, and the differential of $dt\,xL$ is its negative.
The extra terms commute with $x^{-1}$, so the required commutator is still $\hbar I$.
Thus these formulas define an algebra map on every word.
At $t=0$ and $t=1$ they give the two local lifts of the transition.

To retain multiplication on global complexes, an overlap must be represented with its homotopy as well as its endpoints.
For a diagram of parity differential algebras $A_0\to A_{01}\leftarrow A_1$, define
\[
 \mathcal P(A)=
 \{(a_0,a_1,\gamma):
 \gamma\in\Omega^\bullet(k[t])\otimes A_{01},\
 \gamma(0)=a_0|_{01},\ \gamma(1)=a_1|_{01}\}.
\]
Here $\Omega^\bullet(k[t])=k[t]\oplus dt\,k[t]$ with differential $d t=dt$.
Endpoint evaluation sets $dt=0$.
The product is componentwise, with the tensor sign rule on forms.
The total differential preserves the endpoint equations.
This defines an actual unital differential algebra.

\begin{lemma}\label{lem:qc-global-paths}
A strict map between two such diagrams that is a quasi-isomorphism at all three entries induces a quasi-isomorphism of their algebras $\mathcal P$.
For the chart diagram of $\mathcal D_{\hbar,\nu}$, the map from global sections to $\mathcal P(\mathcal D_{\hbar,\nu})$ by constant paths is a quasi-isomorphism.
\end{lemma}
\begin{proof}
Projection from $\mathcal P(A)$ to $A_0\oplus A_1$ is surjective on graded spaces.
A lift uses the path $(1-t)a_0|_{01}+t a_1|_{01}$.
Its kernel consists of forms whose degree-zero part vanishes at both endpoints.
Integration of the $dt$ coefficient identifies this kernel up to homotopy with $A_{01}$ shifted by one odd degree.
The inverse sends $a$ to $dt\,a$.
The homotopy on a $dt$ coefficient $b(t)$ is
\[
 \int_0^t b(u)\,du-t\int_0^1 b(u)\,du.
\]
It vanishes at both endpoints.
Differentiation shows $dh+hd$ is the identity minus the inclusion followed by integration.
All integrals are polynomial antiderivatives over the characteristic-zero field.

The resulting short exact sequence of complexes has kernel equivalent to the shifted overlap complex and quotient $A_0\oplus A_1$.
For a short exact sequence $0\to K\to P\to Q\to0$, lift a cycle $q$ of $Q$ to an element $p$ of $P$.
Its differential belongs to $K$ and is a cycle there.
Changing either representative or lift changes that cycle by a boundary.
This defines a connecting map $\partial:H^j(Q)\to H^{j+1}(K)$.
Its kernel consists precisely of classes admitting cycle lifts.
A cycle of $K$ becomes a boundary in $P$ precisely when its class belongs to the image of $\partial$.
These observations prove the exact cohomology sequence directly.

Suppose a map between two short exact sequences induces isomorphisms on the kernel and quotient cohomology.
Given a middle class in the target, lift its quotient class through the quotient isomorphism.
Its connecting class vanishes by the kernel isomorphism, so it lifts to a middle class in the source.
The difference from the desired target class comes from kernel cohomology and can be lifted there.
This proves surjectivity on middle cohomology.
For injectivity, a source class mapping to zero first has zero quotient class.
Write it as the image of a kernel class.
The target kernel class lies in the image of the connecting map.
Lift that quotient class, subtract its connecting image, and use injectivity on kernel cohomology.
The original middle class is therefore zero.
Apply this argument to the path-algebra short exact sequences to obtain the first assertion.

For the operator sheaf, each chart differential is zero.
Integration therefore identifies the odd cohomology of $\mathcal P$ with the cokernel of the overlap difference.
Its even cohomology is the kernel of that difference, which is the algebra of global sections.
Lemma~\ref{lem:qc-laurent-splitting} makes the cokernel zero.
Constant paths give the asserted quasi-isomorphism.
\end{proof}

The affine moment complex also maps to the glued normal complex.
Write $C_M=\mathsf W^0\otimes\Lambda(\epsilon)$ with $d\epsilon=M-\hbar\nu$, and reduce its grading to parity.
Localization on the first chart sends
\[
 X_{12}\mapsto v,\quad X_{21}\mapsto x\mu-x^2v,
 \quad X_{11}\mapsto\mu-xv,\quad X_{22}\mapsto xv.
\]
These are the actual localization maps constructed above, with $M$ renamed $\mu$.
They agree on the overlap, and sending $\epsilon$ to the global generator of $\mathcal K$ preserves the differential.
Thus they give $C_M\to\mathcal P(\mathcal K)$ by constant paths.

\begin{theorem}\label{thm:qc-global-normal}
Write $Q=\Gamma(\mathbb P^1,\mathcal D_{\hbar,\nu})$, identified with the central quotient in Theorem~\ref{thm:qc-global-algebra}.
The following diagram consists of unital parity differential-algebra maps, and its square commutes:
\[
 \begin{tikzcd}[column sep=large,row sep=large]
 C_M\arrow[r]\arrow[d]&\mathcal P(\mathcal K)\arrow[d]\arrow[r]&\mathcal P(\mathcal E)\\
 Q\arrow[r]&\mathcal P(\mathcal D_{\hbar,\nu})&
 \end{tikzcd}
\]
Every arrow is a quasi-isomorphism.
Proposition~\ref{prop:qc-bar-contraction} also identifies $C_M$ with the relative moment bar through explicit multiplicative quasi-isomorphisms.
\end{theorem}
\begin{proof}
Apply Lemma~\ref{lem:qc-global-paths} to both strict arrows in Proposition~\ref{prop:qc-normal-sheaf}.
The constant-path arrow is covered by the second assertion of the lemma.
The augmentation $C_M\to Q$ is a quasi-isomorphism by the injectivity of the moment relation and Theorem~\ref{thm:qc-global-algebra}.
The two routes around the square send $\epsilon$ to zero and each $X_{ij}$ to its reduced local operator.
They are algebra maps, and these elements generate $C_M$, so the square commutes on every element.
On cohomology, the top arrow followed by the middle vertical arrow is an isomorphism.
The middle vertical arrow is itself an isomorphism on cohomology, hence so is the top arrow.
All maps preserve multiplication by their displayed construction.
\end{proof}

This construction retains the normal potential and its overlap data while recovering the quantum operator algebra.
Its normal odd generator has square zero.
An odd primitive with nonzero square would require a different comparison, as its product cannot be discarded by these maps.

\section{Period reflection and finite-dimensional states}

The assignments $a_i\mapsto p_i$ and $p_i\mapsto-a_i$ define an automorphism of the Weyl algebra.
They preserve the commutators and send $M$ to $-M-2\hbar$.
Thus they give the period reflection
\[
 \mathsf A_{\hbar,\nu}\cong\mathsf A_{\hbar,-\nu-2},
 \qquad E\mapsto-F,\quad F\mapsto-E,\quad H\mapsto-H.
\]
The meaning of the right side is the same presentation after the displayed substitution for the parameter.
The scalar $\nu(\nu+2)$ is unchanged.

There is a separate reflection for the sheaf.
An opposite algebra has the same elements but reverses every product.
On each chart, the assignments $x\mapsto x$, $v\mapsto-v$ define an anti-homomorphism, meaning a map that reverses multiplication.
Applied to the right side of \eqref{eq:qc-operator-transition}, it gives
\[
 vx^2+\hbar\nu x=x^2v+\hbar(\nu+2)x.
\]
This is the negative of the transition at period $-\nu-2$.
The local anti-homomorphisms consequently glue and prove
\[
 \mathcal D_{\hbar,\nu}^{\mathrm{op}}\cong\mathcal D_{\hbar,-\nu-2}.
\]
The shift also appears in differentials: $dx'=-x^{-2}dx$, so the cotangent line is $\mathcal O(-2)$.

Set $\hbar=1$ and let $\nu=n$ be a nonnegative integer.
Global sections of $\mathcal O(n)$ are the polynomials $k[x]_{\le n}$.
Indeed, a polynomial in the first frame becomes $x'^{\,n}f(x'^{-1})$ in the second, which is regular exactly when $\deg f\le n$.
On the basis $1,x,\ldots,x^n$, the three operators are
\[
 E x^j=jx^{j-1},\qquad
 Hx^j=(n-2j)x^j,\qquad
 Fx^j=(n-j)x^{j+1}.
\]
The last coefficient vanishes at $j=n$.
Thus the subspace is preserved by the actual differential operators, and their relations hold without a projection onto that subspace.
Their central scalar is $n(n+2)$.

For $n=1$, this scalar is $3$.
On two copies of this two-dimensional space, let a generator act by its sum on the two factors.
The alternating vector $1\otimes x-x\otimes1$ is annihilated by $E,F,H$, so its total Casimir is zero.
The symmetric subspace is generated from $1\otimes1$ by $F$.
This highest vector is killed by $E$ and has $H$-eigenvalue $2$, so its Casimir is $2^2+2\cdot2=8$.
Centrality gives the same value throughout that subspace.
Hence the sum action does not preserve the fixed relation $\Omega=3$ on tensor products.
Tensor operations must retain the larger algebra before fixing this central value.

The operator comparison has recovered products, local transitions, a normal differential, and global states with the same shifted central scalar.
A chiral extension must additionally retain the singular products of fields and their derivatives.
The coordinate transition \eqref{eq:qc-operator-transition} determines its required zero-mode image, but does not determine those additional operations.
